% !TeX root = ../main.tex
% ----------------------------------------------------------------------
% Entornos: definiciones, teoremas, ejemplos, observaciones y demostraciones
% ----------------------------------------------------------------------
\usepackage{thmtools}
\usepackage[framemethod=default]{mdframed}

\mdfsetup{skipabove=0.9em, skipbelow=0.6em}

% Caja con barra lateral de color
\mdfdefinestyle{barra}{
  topline=false, rightline=false, bottomline=false,
  linewidth=2.5pt,
  innerleftmargin=10pt, innerrightmargin=8pt,
  innertopmargin=10pt, innerbottommargin=10pt,
  leftmargin=0pt, rightmargin=0pt,
  splittopskip=0.5em, splitbottomskip=0.5em
}

% Teoremas, proposiciones, lemas y corolarios
\declaretheoremstyle[
  headfont=\bfseries\color{darkblue},
  notefont=\bfseries\color{accentblue},
  notebraces={(}{)},
  bodyfont=\normalfont,
  headpunct={.},
  postheadspace=0.5em,
  mdframed={style=barra, linecolor=darkblue, backgroundcolor=darkblue!5}
]{caja-teorema}

% Definiciones
\declaretheoremstyle[
  headfont=\bfseries\color{accentblue},
  notefont=\bfseries\color{accentblue},
  notebraces={(}{)},
  bodyfont=\normalfont,
  headpunct={.},
  postheadspace=0.5em,
  mdframed={style=barra, linecolor=accentblue, backgroundcolor=accentblue!5}
]{caja-definicion}

% Ejemplos
\declaretheoremstyle[
  headfont=\bfseries\color{mediumgray},
  notefont=\bfseries\color{mediumgray},
  notebraces={(}{)},
  bodyfont=\normalfont,
  headpunct={.},
  postheadspace=0.5em,
  mdframed={style=barra, linecolor=mediumgray!60, backgroundcolor=white}
]{caja-ejemplo}

% Observaciones
\declaretheoremstyle[
  headfont=\bfseries\color{mediumgray},
  notefont=\bfseries\color{mediumgray},
  notebraces={(}{)},
  bodyfont=\normalfont,
  headpunct={.},
  postheadspace=0.5em,
  mdframed={style=barra, linecolor=mediumgray!60, backgroundcolor=lightgray}
]{caja-observacion}

\declaretheorem[numberwithin=chapter, style=caja-definicion, name=Definición]{definition}
\declaretheorem[sibling=definition, style=caja-teorema, name=Proposición]{prop}
\declaretheorem[sibling=definition, style=caja-teorema, name=Teorema]{theorem}
\declaretheorem[sibling=definition, style=caja-teorema, name=Lema]{lemma}
\declaretheorem[sibling=definition, style=caja-teorema, name=Corolario]{corollary}
\declaretheorem[sibling=definition, style=caja-ejemplo, numbered=no, name=Ejemplo]{eg}
\declaretheorem[style=caja-observacion, numbered=no, name=Observación]{note}
\declaretheorem[style=caja-observacion, numbered=no, name=Ejercicio]{ejercicio}

% Demostraciones (fuera de las cajas)
\renewcommand{\proofname}{Demostración}
\renewcommand{\qedsymbol}{\(\blacksquare\)}
\makeatletter
\renewenvironment{proof}[1][\proofname]{%
  \par\pushQED{\qed}\normalfont\topsep=0.3em\relax
  \trivlist\item[\hskip\labelsep\itshape\textcolor{darkblue}{#1}\@addpunct{.}]\ignorespaces
}{%
  \popQED\endtrivlist\@endpefalse
}
\makeatother
